\chapter{Introduction}

\begin{sloppypar}The VNB Proof Checker is an interactive proof assistant implemented in
MIT~Scheme.  Its logical foundation is a first-order set theory in the
style of von~Neumann--Bernays (VNB): a single-sorted universe of
\emph{classes}, with \textsc{set} as a distinguished predicate and
membership ($\in$) as the only non-logical primitive.\end{sloppypar}
[[NOTE: References needed. NO HALLUCINATIONS please]]

The proof architecture follows the sequent-calculus approach of the
IMPS~2.0 system~\cite{IMPS}: proofs are \emph{deduction graphs} whose
nodes are sequents and whose edges are inference steps.  Every
inference step is validated by a small, auditable \emph{primitive
inference kernel}; higher-level tactics and a form of conditional
rewrite rules, called macetes in the IMPS system, are built on top of
the inference kernel and cannot introduce unsoundness.
The deduction-graph proof structure and the use of local contexts in
macete application are due to L.~G.~Monk~\cite{Monk1988}, whose work
was a direct influence on IMPS.
[[Reference to IMPS and an online Portuguese language dictionary
(e.g.\ \url{https://dicionario.acad-ciencias.pt/pesquisa/macete/} has
an entry)]]

This project is, among other things, a proof of concept: a rigorous
mathematical proof system built entirely by AI
[[Note: we can reference documentation for Claude]],
with architectural direction from F.~J.~Thayer.  It is inspired by
ideas from IMPS~\cite{IMPS}.

%% =========================================================
\section{VNB: primitives and reductions}
\label{sec:primitives}

\emph{[References to VNB in the published literature will be added here
after a literature search.  No content will be added without accurate
citations.]}

Pure VNB set theory is a spare formalism.  The natural numbers, the
reals, functions, and ordinals can all be \emph{constructed} from the
axioms of set theory---but those constructions are long, intricate, and
bear little resemblance to how mathematics is actually written or thought
about.  A proof that $\sin(\pi/2) = 1$ should not need to unfold the
Dedekind-cut definition of the reals.

This system therefore starts from a broad primitive base.  VNB set
theory supplies the logical foundation, but many mathematical objects
are taken as \emph{primitive constants} and \emph{primitive
constructors}, governed by axioms rather than by explicit
set-theoretic construction:

\begin{itemize}[label = $\circ$, leftmargin= 0.9cm]
\item \textbf{Number systems} \texttt{NN}, \texttt{ZZ}, \texttt{QQ},
  \texttt{RR}, \texttt{CC} are primitive constants denoting sets.  Their
  elements and operations are characterised by ring, field, and order
  axioms, not by Dedekind cuts, Cauchy sequences, or ordinal arithmetic.
\item \textbf{Ordinals} \texttt{ORD} form a primitive proper class with
  its own ordering predicate \texttt{<=\_ORD}, independent of the
  set-membership relation.
\item \textbf{Function spaces} \texttt{fun($A$, $B$)} are primitive
  set constructors.  Functions are not encoded as sets of ordered
  pairs; application \texttt{$f$($x$)} is a primitive operation.
\item \textbf{Mathematical structures} such as groups, vector spaces,
  and Banach spaces are declared with \texttt{def-structure}, which
  installs accessor operations and a typing predicate.
\end{itemize}

In principle every one of these primitives could be reduced to the bare
VNB axioms---ordinals by the von~Neumann construction, reals by Dedekind
cuts, functions by sets of ordered pairs, and so on.  This approach
simply takes these reductions as \emph{existence theorems} that
have been proved (or could be proved) elsewhere, and works with the
objects directly at their natural level of abstraction.  The resulting
proofs are shorter, more readable, and correspond more closely to
informal mathematical practice.

A further departure from standard set theory is the treatment of
\emph{partial application}.  Any syntactic expression \texttt{f(x)}
is a well-formed term, but it may be \emph{undefined} if \texttt{f} is
not a function applicable to \texttt{x}.  The system provides both
strict equality \texttt{a = b} (false if either side is undefined) and
quasi-equality \texttt{a == b} (true when both sides are undefined in
the same way), following the approach of IMPS~\cite{IMPS}.

A consequence of this broad primitive base is that membership in a
primitive constant has \emph{indeterminate truth value} in general.
The formula \texttt{x in 7} is syntactically well-formed --- VNB is
single-sorted and \texttt{in} accepts any two class expressions ---
but the system makes no claim about whether it holds. Please note that
this does not mean that the formula is undefined: whether it is true
or not depends on how $7$ is interpreted in VNB. This system
takes $7$ as a primitive number constant and makes no identification
of it with any set-theoretic object.  Membership statements of this
kind are grammatically meaningful but proof-theoretically inert.

%% =========================================================
\section{Architecture}
\label{sec:architecture}

\subsection{Proof architecture}

\begin{sloppypar}
The proof architecture follows the sequent-calculus approach of the
IMPS~2.0 system~\cite{IMPS}: proofs are \emph{deduction graphs} whose
nodes are sequents and whose edges are inference steps.  Every inference
step is validated by a small, auditable \emph{primitive inference kernel};
higher-level tactics and macetes are built on top of it and cannot
introduce unsoundness.
\end{sloppypar}

\subsection{Theory architecture}

In contrast to the \emph{little theories} methodology of
IMPS~\cite{IMPS}---where each mathematical domain inhabits its own
separate mini-theory---VNB works with a \emph{single, growing current
theory}.  Mathematical structures are declared directly in that one theory
using \texttt{def-structure}, and their axioms accumulate there alongside
the base VNB axioms and the primitive number-system axioms.  This reflects
the view of mathematics as one unified body of knowledge rather than a
federation of isolated theories.

The system does permit \emph{minimal} or even \emph{odd-ball theories}
for the cases where a genuinely separate context becomes necessary.
These theories can be created with \texttt{make-theory} and carry
the primitive class constants but \emph{no structural axioms}.  Such
theories are expected to be exceptional; all routine mathematical work
takes place in the single current theory.

\begin{remark}\label{remark:27137-61828-986360}
%
For the potential user interested in \emph{constructivist} mathematics,
this system is not appropriate. These potential users should consider
alternatives such as Coq, which has a very well established library of
mathematics.
%
\end{remark}

\subsection{Design goals}

\begin{itemize}[ label = $\circ$, leftmargin = 1cm]
\item \textbf{Enriched set-theoretic foundation.}  VNB set theory with
  primitive constants for standard mathematical objects; no type
  hierarchy; partial function application with undefinedness.
\item \textbf{Transparent kernel.}  All proof steps reduce to a small
number of
  primitive inferences; the kernel is under 300 lines of Scheme.
\item \textbf{Composable rewriting.}  A macete mechanism (borrowed from
  IMPS) allows theorems to be automatically converted into named rewrite
  strategies that can be sequenced and combined.
\item \textbf{N-ary constructors.}  Ordered tuples \texttt{[$e_1$,
  \ldots, $e_n$]} and Cartesian products \texttt{cartesian($A_1$, \ldots,
  $A_n$)} are first-class, supporting the definition of mathematical
  structures with multiple carriers and operations.
\end{itemize}

\subsection{Why MIT Scheme?}

MIT~Scheme was chosen for several practical reasons: it is small (the
proof kernel is under 300 lines of idiomatic Scheme), easy to learn
and write code, free and readily available on all major platforms, and has
a long track record of correctness and stability.  Since all code is
open source, porting to another language with AI assistance should be
straightforward.

%% =========================================================

%% =========================================================
\section{Quick start}
\label{sec:quickstart-chap}

\subsection{Installation and startup}

\subsubsection{Requirements}

MIT~Scheme~11.2 or later, and for best interactive use one of
GNU/Emacs or XEmacs.  With minor modifications, the native Scheme text
editor Edwin also works.  No other dependencies.

\subsubsection{Directory layout}

\begin{VNBsnip}
~/prover/
  prover                 executable script
  load.scm               loader (sources all files in order)
  expressions.scm        expression language
  wff.scm                <wff> record type, wff-child, wff-in-theory, wff-equiv?
  sequents.scm           sequents over <wff>; context operations
  deduction-graphs.scm   proof graph data structures
  primitive-inferences.scm  inference kernel
  macetes.scm            macete mechanism
  theory.scm             theory management + VNB base axioms
  structures.scm         mathematical structure definitions
  number-systems.scm     NN, ZZ, QQ, RR, CC constants and axioms
  contexts.scm           local context management, make-wff, unravel
  arith-eval.scm         ground arithmetic decision procedure
  proof-commands.scm     high-level proof commands
  interactive.scm        interactive short-form commands (sp, di, ai, ...)
  vnb.el                 GNU Emacs / XEmacs interface
\end{VNBsnip}

\subsection{Starting the checker}

\begin{VNBsnip}
~/prover/prover              # interactive REPL
~/prover/prover proof.scm    # run a proof script and exit
~/prover/prover -i proof.scm # run a script, then stay interactive
\end{VNBsnip}

The interactive session is a standard MIT~Scheme REPL with the proof
checker pre-loaded.
\subsection{The Emacs interface}
\label{sec:emacs}

The file \texttt{vnb.el} provides a GNU~Emacs and XEmacs interface for
the proof checker.  It requires no external packages beyond
\texttt{comint}, which ships with every Emacs distribution.

\subsubsection{Installation}

Add the following to your Emacs init file (\texttt{\textasciitilde/.emacs}
or \texttt{\textasciitilde/.xemacs/init.el}):

\begin{verbatim}
(add-to-list 'load-path "~/prover")
(require 'vnb)
\end{verbatim}

\noindent
\textbf{Do not byte-compile \texttt{vnb.el}.}  GNU~Emacs and XEmacs use
incompatible byte-code formats; loading the source file interpreted is
the correct approach for both.

\subsubsection{Starting the interface}

\begin{VNBsnip}
M-x vnb
\end{VNBsnip}

This starts the prover subprocess (if not already running), and opens a
two-buffer layout:
\begin{itemize}[label=\textendash, leftmargin=1cm]
\item \textbf{Left window --- \texttt{*VNB*} REPL.}  A \texttt{comint}
  buffer running MIT~Scheme with the proof checker pre-loaded.  Type any
  Scheme expression and press \texttt{RET} to evaluate it.
\item \textbf{Right window --- \texttt{*VNB State*}.}  A read-only
  display buffer that is updated automatically after every command that
  calls \texttt{(show)}.  It always shows the current proof state: open
  goals, assumptions, and focus.
\end{itemize}

The two buffers are kept synchronised by a comint output filter that
watches for the sentinel lines \texttt{;;VNB-STATE-BEGIN} /
\texttt{;;VNB-STATE-END} emitted by \texttt{(show)} and routes the
enclosed text to the state buffer.

\subsubsection{Key bindings in the \texttt{*VNB*} REPL buffer}

\begin{itemize}[label = $\circ$, leftmargin=1cm]
\item \texttt{RET} Submit expression to the prover (standard comint).
\item \texttt{C-c C-s} Send \texttt{(show)} to refresh the state display.
\item \texttt{C-c C-p} Switch to the \texttt{*VNB State*} buffer.
\end{itemize}

\subsubsection{Key bindings in the \texttt{*VNB State*} buffer}

\begin{itemize}[label = $\circ$, leftmargin=1cm]
\item \texttt{q} Return to the \texttt{*VNB*} REPL buffer.
\item \texttt{C-c C-s} Send \texttt{(show)} to refresh the state display.
\end{itemize}

\subsubsection{The \texttt{*VNB Commands*} scratch buffer}

\begin{VNBsnip}
M-x vnb-command-buffer
\end{VNBsnip}

Opens a \texttt{*VNB Commands*} buffer in \texttt{vnb-command-mode}.
This buffer works exactly like the Emacs \texttt{*scratch*} buffer (which
uses \texttt{lisp-interaction-mode}), except that \texttt{C-j} sends
expressions to the \emph{VNB prover} instead of evaluating them as Emacs
Lisp.

Write proof commands as S-expressions in the buffer.  Position point
immediately after any complete expression and press \texttt{C-j}:
the expression is sent to the running VNB process, and the return value
is inserted on a new line below point.  The buffer accumulates a
readable session transcript.

\begin{VNBcode}
(sp (make-wff-from-string "n in NN and m in NN implies m in NN and n in NN"))
;;VNB-STATE-BEGIN
1 open goal(s).
Focus:
  [0] ()  =>  (((n in nn) and (m in nn)) implies ((m in nn) and (n in nn)))
;;VNB-STATE-END

(di)       C-j
;;VNB-STATE-BEGIN
2 open goal(s).
Focus:
  [1] ((n in nn) and (m in nn))  =>  ((m in nn) and (n in nn))
Other open goals:
  [0] ()  =>  (((n in nn) and (m in nn)) implies ((m in nn) and (n in nn)))
;;VNB-STATE-END
\end{VNBcode}

\paragraph{Key bindings in \texttt{vnb-command-mode}.}

\begin{itemize}[label = $\circ$, leftmargin=1cm]
\item \texttt{C-j} Send S-expression before point to VNB; insert result below.
\item \texttt{C-c C-c} Send entire buffer to VNB; append result.
\item \texttt{TAB} Complete VNB command name at point.
\item \texttt{C-c C-d} Describe VNB command: show signature and documentation.
\item \texttt{C-c C-s} Refresh \texttt{*VNB State*} proof-state buffer.
\end{itemize}

\subsubsection{Command reference and completion}

\begin{sloppypar}\texttt{vnb.el} maintains the list \texttt{vnb-commands} of all proof
commands, each with a usage signature and a one-line description.  The
derived alist \texttt{vnb-commands-alist} maps short names to signatures
and is used by \texttt{TAB} completion in \texttt{vnb-command-mode}.\end{sloppypar}

\begin{VNBsnip}
M-x vnb-describe-command   RET   di   RET
; Echo area: (di)
;   Direct inference: decompose current goal by its top connective. ...
\end{VNBsnip}

\subsubsection{Programmatic evaluation: \texttt{vnb-eval-string}}

For scripting or integration with other Emacs packages:

\begin{VNBsnip}
; str is a Scheme expression as a string
(vnb-eval-string str)           ; => Scheme return value as string
(vnb-eval-string str timeout)   ; => same, with explicit timeout (seconds)
\end{VNBsnip}

\texttt{(vnb-eval-string STR \&optional TIMEOUT)} sends the string
\texttt{STR} to the running prover, waits up to \texttt{TIMEOUT} seconds
(default 30) for the next REPL prompt, and returns the Scheme return
value as an Emacs string.  It signals an error if the prover is not
running.

%% =========================================================
\subsection{A complete proof}
\label{sec:quickstart}

\begin{sloppypar}
After installing the Emacs interface (Section~\ref{sec:emacs}), launch
the prover with \texttt{M-x vnb}.  Two buffers appear: \texttt{*VNB*}
(the REPL) and \texttt{*VNB State*} (the live proof state).  Type the
following five commands at the \texttt{*VNB*} prompt, pressing
\texttt{C-j} (or RET) after each:
\end{sloppypar}

\begin{VNBsnip}
(sp (make-wff-from-string "n in NN and m in NN implies m in NN and n in NN"))
(di)                                   ; assume both, prove (m in nn) and (n in nn)
(ai "n in NN and m in NN")            ; split into (n in nn) and (m in nn)
(di)                                   ; split goal: prove (m in nn), then (n in nn)
(ass) (ass)                            ; both are in the assumption list
\end{VNBsnip}

The \texttt{*VNB State*} buffer updates after each command.  The final
state shows \texttt{Proof complete.}

\paragraph{What just happened.}
\texttt{(make-wff-from-string \ldots)} parses a string in the VNB string
syntax and wraps the result as a typed \texttt{wff} in the current
theory.  \texttt{(sp \ldots)} starts a new proof.  \texttt{(di)} is
direct-inference: it decomposes the goal by its outermost connective.
\texttt{(ai \ldots)} is antecedent-inference: it decomposes a named
assumption; the argument is a formula string (or a raw S-expression for
low-level use) that identifies the assumption by alpha-equivalence.
\texttt{(ass)} discharges a goal that already appears in the assumption
list.

The same proof at the Scheme level (no short forms) reads:
\begin{VNBsnip}
(let* ((wic (make-wff-from-string
              "n in NN and m in NN implies m in NN and n in NN"))
       (ps  (start-proof wic))
       (ps  (cmd-direct-inference ps))
       (ps  (cmd-antecedent-inference ps (parse-string "n in NN and m in NN")))
       (ps  (cmd-direct-inference ps))
       (ps  (cmd-assumption ps))
       (ps  (cmd-assumption ps)))
  (print-proof-state ps))
\end{VNBsnip}
The interactive short forms in \texttt{interactive.scm} are simple
wrappers: each mutates a global \texttt{*ps*} and calls \texttt{(show)}
to reprint the state.  All examples in this manual use the short forms;
the long form is available when scripting non-interactively.

\paragraph{What \texttt{make-wff} checks.}
\texttt{make-wff} (and \texttt{make-wff-from-string}, which calls it)
performs a position-typed walk of its argument and signals an error
rather than wrap a malformed expression.  The checks are:
\begin{sloppypar}
\begin{itemize}[label=$\ast$, leftmargin=1cm]
\item \textbf{Numbers in wff positions} are rejected:
  \texttt{(make-wff-from-string "1 and 35")} $\Rightarrow$ ``number in wff
  position 1''.
\item \textbf{Term-forming heads in wff positions} are rejected:
  \texttt{(make-wff-from-string "choice(A)")} as a top-level wff fails because
  \texttt{choice} produces a term, not a wff.
\item \textbf{Wff-forming heads in term positions} are rejected:
  \texttt{(make-wff-from-string "(A and B) in C")} fails because \texttt{A and B}
  is a wff and the first argument of \texttt{in} must be a term.
\item \textbf{Bare symbols in wff positions} are rejected.  Only
  \texttt{truth} and \texttt{falsity} are valid bare-symbol wffs;
  anything else (e.g.\ \texttt{A implies B}) is an error.
\item \textbf{Free-symbol role conflicts}: a symbol that appears free
  in both wff and term positions in the same formula is rejected.
\item \textbf{Wff-position binders}: a binder whose body uses the
  bound variable in wff position is rejected within its scope.  This
  catches \texttt{forall([p], p implies p)}: \texttt{p} is bound as a
  class variable but the body uses it as a proposition.
\item \textbf{Bound vs.\ free same-name use} is \emph{accepted with a
  warning}.  In \texttt{forall([a], a in nn) implies a in zz} the
  inner \texttt{a} is bound and the outer \texttt{a} is free.  The
  validator emits a \texttt{;; warning} but does not reject.
\item \textbf{Arities and bound-variable shapes} are checked: every
  connective and quantifier gets the right number of arguments;
  \texttt{forall}, \texttt{forsome}, \texttt{iota},
  \texttt{lambda}, \texttt{lambdoid}, and \texttt{sep} bound variables must be
  symbols.
\end{itemize}
\end{sloppypar}
%
\begin{remark}\label{remark:27135-34761-138906}
%
The set-theoretic \texttt{power(A)} (power set) and arithmetic
\texttt{x \^{} n} (exponentiation) share the same underlying symbol because
MIT~Scheme reads identifiers case-insensitively; \texttt{make-wff}
distinguishes them by arity (1 vs.\ 2 arguments).  All downstream
operations (\texttt{free-vars}, \texttt{subst-free},
\texttt{alpha-equiv?}, macete pattern matching) honour the same
arity-based dispatch.

%
\end{remark}

%% =========================================================
\subsection{A macete example}
\label{sec:quickstart-macete}

Macetes are named rewrite rules derived from installed theorems.
To apply a macete outside a proof, install a theorem and call
\texttt{apply-macete}:

\begin{VNBcode}
(install-theorem! 'eq-symm
  (make-wff-from-string "forall([x, y], x = y iff y = x)"))

(apply-macete 'eq-symm (make-wff-from-string "3 + 5 = 8"))
; => wff for "8 = 3 + 5"

(apply-macete 'eq-symm (make-wff-from-string "n = m implies truth"))
; => wff for "m = n implies truth"

(apply-macete 'eq-symm (make-wff-from-string "3 in nn"))
; => #f          -- pattern x = y does not match; no rewrite fires
\end{VNBcode}

\texttt{apply-macete} accepts and returns a \texttt{wff} object; it
returns \texttt{\#f} when the pattern matches nowhere.
To apply a macete as a proof step, use
\texttt{(mac \textit{name})} inside a proof (see
Section~\ref{sec:macetes}).

%% =========================================================


%% =========================================================
